\documentclass[11pt]{article}
\usepackage{amsmath,hyperref}
\title{Prepared and Armed Successors Have Different Cancellation Cuts}
\date{2026-11-27}
\begin{document}
\maketitle

\paragraph{Audit.}
The original writer has already been fenced and its simulator
checkpoint made available to two prospective successors.
There is one named unsettled warehouse intent. A separate trusted,
durable, single SQLite registry can choose an intended successor.
Each successor has its own native source state and can remain
executable despite being unreachable from the selector. The parties
are managed, compliant processes. Trusted logical executor origins,
paths, and grant/closure provenance are supplied by this fixture.
The registry is not a consensus implementation; the typed proof
values are in-process data, not deployed credentials. The effect is
the authoritative simulator step in the selected destination
database, not an arbitrary external actuator. Times below describe
ordering, not measured wall-clock deadlines.

\paragraph{Admission cut.}
At most one update of a single origin choice wins each SQLite registry
transaction. A prepared nomination is not permission to install;
the current nominee and generation must successfully arm it first.
Before arming, selection can be changed without reaching the prior
candidate, because the atomic registry admits either the old arm
or the selection update, not both. After the arm transaction commits,
even a lost reply does not make that authority unused. An old installer
may already be acting on its earlier read of the arm. Rewriting the
registry row alone does not touch its local mutation path.

Reassignment after arm therefore first closes the old nominee at
its effect store. The local closure is durable and serialized against
both installation and native publication: an earlier outstanding arm
cannot turn a blocked store back into a writable one. If the nominee
had installed, closure also fences its active publication epoch
and seals its complete named native state. Transfer from this
actual closed store rather than the original donor's image.
The genuine receipt controls whether the next worker applies once or
returns the stored result. If the nominee had not installed,
the local block permits reuse of the original frozen snapshot for
the different destination. The already closed local gate is
not copied as the new candidate's executor identity.

\paragraph{Checking the boundary.}
A nominally valid donor seal accepted twice is the previous
\href{https://github.com/eternal-november/eos-research-programmes/pull/742}
{two-copy control}. A write-once recipient row remedies that fork
for initial choice: four claim crashes and four subprocess races
select one immutable recipient and one native application each.
It cannot simply be overwritten when that recipient is already
writable. This write-once comparison is a diagnostic control;
the established comparison is a real fencing mechanism.

The new sweep crosses five old-candidate phases with actual local
store access or withheld access. Before arm the new candidate can
apply without access to the old. After arm, missing closure withholds
the dependent request while a separate native carrier moves. A
committed local closure enables replacement even after an old arm
was withheld; if a native effect committed before closure, the new
leader returns its named receipt. If only staged, the old subprocess
exit rolls back and the successor instead applies once. There are
two actual registry crashes bracketing arm commit and four independent
subprocess races between arm and reassignment. The race wins once,
with either accepted order. A held native writer prevents the local
closure transaction from starting; when its effect commits first,
the copied receipt is observed on the successor. Omitting the
effect-local barrier lets the old paused installer and replacement
both apply, despite the registry choosing only one current candidate.

\paragraph{Independent safety projection.}
For the one-intent family, selection of B implies either that no
arm went to A or that A is locally closed. An open A implies an
issued arm. A copied state at B must have been selected and retain
whether A paid at closure. Since A cannot newly pay after this
selection, B's copied paid bit remains faithful. B can newly apply
only when this bit was false. Lean checks the invariant under
abstract steps; the independent finite enumerator finds 37 reachable
states and 48 transitions in its exact finite kernel, no double pay.
Dropping the arm gate, dropping the local gate, or copying the
pre-nominee state after a committed effect each gives a shortest
duplicate-effect trace. The finite projection does not prove the
database/runtime refinement, multiple in-flight intents, real
quiescence, or a notification schedule.

\paragraph{Alternative and residual.}
A resource that validates an up-to-date sequencer on its own
mutation path can implement an ordinary direct fencing design.
This is the relevant reference, not a registry that only overwrites
an appointment, and the distinction is already present in
\href{https://www.usenix.org/legacy/event/osdi06/tech/full_papers/burrows/burrows_html/}
{Chubby's lock sequencers}. Where the two destinations can instead
mutate the very same authority-enforced effect store, there is no
need to transfer two volatile writable copies. Neither construction
is credited if its effect-store or timing assumptions cannot be met.
The new result is a checked choice of boundary and a backed-up
counterexample, not an advantage over a matched direct fencing
design. An unreachable armed installer whose scope lacks an
effect-side fencing channel still prevents this recovery by the
wished-for deadline. Keeping independent work moving is
containment, not completion of the blocked task. This is held-in
evidence, not a contract acceptance or a field claim.
\end{document}
